% ============================================================================
% Seção: Desenvolvimento orientado por IA
% Marcadores [VERIFICAR] indicam afirmações ainda sem fonte confirmada.
% Os comentários "% NOTA:" são para os autores e devem ser removidos.
% Pacotes necessários: minted, tabularx, longtable, float (macros \codigoarq
% e contador contcodigo definidos no preâmbulo, como já combinado).
% ============================================================================


Esta seção descreve o processo de desenvolvimento assistido por agentes de IA proposto para a implementação do Seed no TCC2. O processo combina especificações explícitas, agentes com permissões restritas, desenvolvimento orientado a testes e verificações de segurança, mantendo as decisões de arquitetura e de requisitos sob responsabilidade do engenheiro. Trata-se de uma proposta de método, cuja efetividade será avaliada no TCC2 (Subseção~\ref{subsec:avaliacao-ia}), e não de uma garantia de qualidade.

%----------------------------------------------------------------------------
\subsection{Fundamentos e limitações dos modelos de linguagem}
\label{subsec:llm-limites}

Modelos de linguagem de grande porte (\textit{Large Language Models}, LLMs) produzem texto de forma sequencial: a cada passo, escolhem a próxima unidade, denominada \textit{token}, a partir de uma distribuição de probabilidades condicionada ao contexto já existente. Como a escolha envolve amostragem, o mesmo \textit{prompt} pode gerar respostas diferentes, o que torna os resultados pouco padronizados. MATHEWS 2024 fixam a semente e a temperatura em seus experimentos para melhorar a replicabilidade, mas reconhecem que algum grau de imprevisibilidade permanece. (AKITA, 2026).

% EM https://www.youtube.com/watch?v=U3bZavG8qQY
% NOTA: a anotação "cerca de 40% do conteúdo serve para agradar o usuário" e a
% analogia da "lootbox" foram removidas por não terem fonte localizada.

Outra limitação é a tendência de modelos a concordar com a premissa do usuário, fenômeno conhecido em inglês como \textit{sycophancy}. Na prática, perguntar se um texto contém erros e perguntar se ele está perfeito pode produzir respostas de teor oposto. [VERIFICAR: citar Sharma et al.] Essa tendência justifica uma regra do processo proposto: o agente não decide requisitos nem arquitetura, e as ambiguidades são devolvidas ao engenheiro por meio da etapa de esclarecimento descrita na Subseção~\ref{subsec:sdd}.


A qualidade do código gerado também preocupa. \citeonline{pearce_2021} geraram 1.689 programas em 89 cenários associados a fraquezas conhecidas de segurança com o GitHub Copilot, e uma parcela expressiva continha vulnerabilidades potencialmente exploráveis. [VERIFICAR: confirmar o percentual diretamente no artigo.] \citeonline{perry_2022} constataram, em estudo com participantes, que quem usou um assistente de IA escreveu código menos seguro e acreditou mais que ele era seguro. Esses resultados sustentam duas escolhas do processo: o engenheiro permanece como decisor e a verificação de segurança é automatizada e independente do agente. CONCLUIR AQUI


\subsection{\textit{Vibe Coding} x Codificação Agêntica}
\label{subsec:agentic}

Popularizado por Andrej Karpathy em fevereiro de 2025, o termo \textit{vibe coding} refere-se a programar delegando à IA a escrita do código e dispensando a sua leitura. A literatura recente tem insistido que usar um LLM, por si só, não basta para caracterizar a prática. Segundo a fronteira traçada por Simon Willison, o código gerado por um modelo, mas revisado, testado e compreendido, equivale a empregar o LLM como assistente de digitação, e não o \textit{vibe coding} \cite{vibe_review_2026}. O que define a prática, nessa leitura, é o afastamento do desenvolvedor em relação ao código.

A postura oposta é chamada por \citeonline{beck_augmented_2025} de \textit{augmented coding}: mantêm-se os valores da engenharia de software tradicional, como testes, código limpo e \textit{commits} pequenos, e se difere somente quando digitado manualmente. \citeonline{sapkota_2025} separam o \textit{vibe coding}, de caráter interativo e exploratório, da \textit{agentic coding}, em que agentes planejam, executam, testam e iteram com pouca intervenção humana, e sugerem que os dois paradigmas sejam combinados em um ciclo centrado no humano.

Segundo \citeonline{ge_vibe_2025}, as práticas de desenvolvimento com agentes de código podem ser agrupadas em cinco modelos: (I) automação irrestrita, em que o agente atua com mínima restrição; (II) colaboração conversacional iterativa, baseada no diálogo contínuo entre desenvolvedor e agente; (III) desenvolvimento orientado a planejamento, em que o trabalho parte de um plano previamente definido; (IV) desenvolvimento orientado a testes; e (V) desenvolvimento aprimorado por contexto, no qual se fornece ao agente informação estruturada sobre o projeto. Sendo assim, o êxito da prática não decorre apenas da capacidade do agente, mas sim, depende da engenharia de contexto, de ambientes de desenvolvimento consolidados e da colaboração entre humano e agente. Nesse cenário é proposto o termo \textbf{codificação agêntica}, que designa o uso de agentes sob práticas do XP: TDD, integração contínua a cada \textit{commit} e refatoração contínua o qual será utilizado para o desenvolvimento do projeto proposto discutido na Subseção~\ref{subsec:subagentes}.


%----------------------------------------------------------------------------
\subsection{Spec-Driven Development}
\label{subsec:sdd}

No desenvolvimento convencional, a especificação costuma orientar a escrita do código e é descartada depois, o desenvolvimento orientado por especificações (\textit{Spec-Driven Development}, SDD) inverte essa lógica \cite{github_speckit}. A especificação é escrita antes de o código ser gerado com IA e passa a ser a fonte de verdade tanto para o humano quanto para o modelo \cite{bockeler_sdd_2025}. Como o termo ainda não tem definição estável, \citeonline{bockeler_sdd_2025} distingue três níveis de adoção: \textit{spec-first}, em que a especificação é escrita antes e guia a tarefa em curso; \textit{spec-anchored}, em que é mantida após a entrega para apoiar a evolução da funcionalidade; e \textit{spec-as-source}, em que é o único artefato editado pelo humano e o código é gerado a partir dela.

O \textit{Spec Kit}, kit aberto do GitHub, conduz o trabalho pelas etapas \textit{specify}, \textit{plan}, \textit{tasks} e \textit{implement}, com etapas opcionais de esclarecimento (\textit{clarify}) e análise \cite{github_speckit}. Os insumos do desenvolvimento são as histórias de usuário e os requisitos funcionais reunidos na especificação (\textit{feature specification}), artefato análogo ao documento de requisitos do produto, comparação que o próprio \citeonline{bockeler_sdd_2025} usa para caracterizar uma especificação. A etapa \textit{clarify} esclarece as lacunas dessa especificação por meio de perguntas ao engenheiro, antes do planejamento \cite{github_speckit}. O plano (\textit{plan}) reúne as decisões técnicas, como tabelas, bibliotecas e contratos de interface, e a lista de tarefas (\textit{tasks}) divide o trabalho em unidades pequenas e revisáveis. Como os artefatos de cada \textit{feature} seguem a mesma estrutura, a revisão fica mais simples e é possível rastrear histórias de usuário, requisitos (FUNC), tarefas e testes. A tabela de rastreabilidade é apresentada no capítulo de resultados.


\citeonline{bockeler_sdd_2025} observa que um fluxo único de SDD pode ser desproporcional para problemas pequenos e que a revisão de muitos arquivos \textit{markdown} é trabalhosa. A autora defende que o controle sobre o que se constrói vem de passos pequenos e iterativos, e não de uma especificação extensa elaborada de antemão. O projeto proposto atende a essa recomendação ao aplicar o fluxo em ciclos curtos, uma história de usuário por vez, de modo que o sistema se forma pela entrega sucessiva de histórias de escopo reduzido, em consonância com o Scrum e o Kanban adotados. Em cada ciclo, a etapa de esclarecimento define os critérios de aceitação e os limites numéricos das tarefas, valores que fundamentam os testes (Subseção~\ref{subsec:subagentes}).

%----------------------------------------------------------------------------
\subsection{Ambiente de agentes}
\label{subsec:ambiente}

\subsubsection{Agentes e subagentes}
\label{subsubsec:agentes}

Para distinguir IA generativa convencional de agente, adota-se a definição da Anthropic, que classifica como sistemas agênticos tanto os \textit{workflows}, em que LLMs e ferramentas são orquestrados por caminhos de código predefinidos, quanto os agentes, em que o LLM dirige dinamicamente o próprio processo e o uso de ferramentas \cite{anthropic_agents}. A capacidade de ler e executar arquivos diretamente no computador, sem copiar e colar no contexto da conversa, decorre do uso de ferramentas, não é a definição em si. A mesma fonte recomenda a solução mais simples possível e observa que a autonomia dos agentes implica custos maiores e risco de acúmulo de erros. [VERIFICAR: confirmar essas ressalvas no texto original.] Por isso, o processo proposto mantém pontos de controle humanos, aproximando-se de um \textit{workflow} supervisionado.

O ganho de acesso traz um problema de segurança: um agente com permissões amplas pode ler arquivos de variáveis de ambiente, acessar dados sensíveis ou sair do diretório do projeto. [VERIFICAR: substituir a referência a material de curso por documentação pública da Anthropic.] Como o Seed trata dados de saúde mental de estudantes, essa preocupação é central e motiva o isolamento descrito na Subseção~\ref{subsubsec:isolamento}.

Neste trabalho, o orquestrador é a sessão principal do agente (Claude Code), que conduz o esclarecimento com o engenheiro e delega tarefas a subagentes. Subagentes possuem contexto próprio e são definidos em arquivos com \textit{frontmatter} YAML delimitado por \texttt{---}: \texttt{name} identifica o subagente; \texttt{description} orienta a delegação; \texttt{model} seleciona o modelo; \texttt{tools} enumera as ferramentas; \texttt{skills} indica habilidades pré-carregadas; e \texttt{hooks} associa verificações ao ciclo de execução \cite{anthropic_subagents}. O desenho não depende de subagentes que criem outros subagentes. [VERIFICAR: confirmar essa possibilidade na documentação.] O Seed define dois subagentes, \texttt{tester} e \texttt{dev}, detalhados na Subseção~\ref{subsec:subagentes}. As ferramentas e os formatos de definição de agentes variam entre fornecedores, e o Spec Kit é concebido para funcionar com vários agentes \cite{github_speckit}, o que reduz a dependência de uma única ferramenta nas especificações, embora não nos arquivos de agente. [VERIFICAR: confirmar a lista de agentes compatíveis.]

% NOTA: o parágrafo sobre "problema energético de 2025" e os exemplos
% "Empresa X/Y" foram removidos: a relação causal entre consumo de energia e o
% surgimento dos agentes não tem fonte e não é sustentada pela literatura localizada.

\subsubsection{Skills}
\label{subsubsec:skills}

Skills são pastas de instruções, descritas em arquivos \texttt{SKILL.md}, que o agente utiliza para tarefas específicas, como revisão de código, mensagens de \textit{commit} ou geração de testes \cite{anthropic_skills}. O agente consulta primeiro a descrição de cada skill e carrega o conteúdo completo apenas quando a tarefa corresponde a ela, o que evita ocupar o contexto com instruções desnecessárias. [VERIFICAR: confirmar esse carregamento progressivo na documentação.] Por isso, a descrição deve informar o que a skill faz e o que provoca seu uso. Skills de leitura e análise podem ter as ferramentas limitadas por meio do campo \texttt{allowed-tools}, reduzindo pedidos de permissão e o risco de alterações indevidas. Skills podem ainda ser associadas a subagentes por meio do campo \texttt{skills}, e a possibilidade de encadeá-las entre si deve ser verificada na documentação [VERIFICAR].

A Tabela~\ref{tab:skills-propostas} lista as skills propostas para o Seed.

\begingroup
\renewcommand{\arraystretch}{1.5}
\begin{table}[H]
    \centering
    \caption{Skills propostas para o processo de desenvolvimento.}
    \label{tab:skills-propostas}

    \begin{tabularx}{\textwidth}{|l|l|X|l|}
        \hline
        \textbf{Skill} & \textbf{Usuário} & \textbf{Função} & \textbf{Situação} \\
        \hline
        \texttt{grill-me} (ou \textit{clarify}) & Orquestrador & Conduzir perguntas ao engenheiro até fechar critérios de aceitação, limites numéricos e decisões técnicas, sem que o agente decida por conta própria. & Conteúdo a descrever [VERIFICAR] \\
        \hline
        \texttt{unit-tests-boundary} & \texttt{tester} & Gerar testes unitários (JUnit~5 ou Jest) com partição de equivalência e valores de borda a partir da especificação. & Declarada \\
        \hline
        \texttt{implement-and-refactor} & \texttt{dev} & Implementar o mínimo para passar nos testes e refatorar com os testes verdes, seguindo as camadas Controller, Service, Model e Repository. & Declarada \\
        \hline
    \end{tabularx}

    \vspace{0.5em}
    \small Fonte: autoria própria.
\end{table}
\endgroup

\subsubsection{Hooks}
\label{subsubsec:hooks}

Hooks são comandos definidos pelo usuário que o Claude Code executa em pontos fixos do ciclo de uso das ferramentas, sem depender de o modelo lembrar de uma instrução \cite{anthropic_hooks}. O evento \texttt{PreToolUse} ocorre antes da chamada de uma ferramenta e permite recusá-la, e \texttt{PostToolUse} ocorre depois que ela é concluída. Um \textit{hook} que termina com código de saída 2 bloqueia a ação, e o motivo retorna ao agente \cite{anthropic_hooks}. Existem ainda eventos como \texttt{Stop}, \texttt{SubagentStop}, \texttt{UserPromptSubmit} e \texttt{PreCompact}, e o conjunto muda entre versões. [VERIFICAR: conferir a lista atual na página oficial.] O critério de uso é objetivo: o que precisa ocorrer sempre é implementado como \textit{hook}, e não como pedido no \textit{prompt}, o que dialoga com a degradação de instruções em contextos longos discutida na Subseção~\ref{subsec:llm-limites}.

Hooks não invocam skills nem comandos de conversa. Portanto, o esclarecimento de requisitos é iniciado pelo orquestrador, por solicitação do engenheiro ou por regra no arquivo \texttt{CLAUDE.md}, que reúne o onboarding do projeto para o agente. Os \textit{hooks} atuam como travas, e não como gatilhos. Um formatador, por exemplo, é executado depois que o agente edita um arquivo, e não antes de executar código.

No Seed, os \textit{hooks} têm duas funções: limitar o acesso do agente a variáveis de ambiente e a conteúdo fora do diretório de código, e aplicar regras de formatação e limpeza a cada geração de código. A Tabela~\ref{tab:hooks-propostos} apresenta o conjunto proposto.

\begingroup
\small
\renewcommand{\arraystretch}{1.4}
\begin{longtable}{|l|l|>{\raggedright\arraybackslash}p{0.42\textwidth}|>{\raggedright\arraybackslash}p{0.17\textwidth}|}
    \caption{Hooks propostos para o processo de desenvolvimento.}\label{tab:hooks-propostos}\\
    \hline
    \textbf{Evento} & \textbf{Seletor} & \textbf{Função} & \textbf{Escopo} \\
    \hline
    \endfirsthead
    \hline
    \textbf{Evento} & \textbf{Seletor} & \textbf{Função} & \textbf{Escopo} \\
    \hline
    \endhead
    PostToolUse & \texttt{Edit|Write} & Formatar o arquivo alterado conforme a linguagem e devolver ao agente os erros de \textit{lint}. As ferramentas serão definidas [VERIFICAR]. & Global \\
    \hline
    PreToolUse & \texttt{Read|Edit|Write|Bash} & Bloquear acesso a arquivos de variáveis de ambiente, chaves e credenciais, como segunda camada após o isolamento do contêiner. & Global \\
    \hline
    PreToolUse & \texttt{Bash} & Bloquear comandos destrutivos, como remoção recursiva, \textit{push} forçado, \textit{reset} rígido e \texttt{DROP}. & Global \\
    \hline
    PreToolUse & \texttt{Read|Grep|Glob|Write|Edit} & Lista de permissões (\textit{allow-list}): o \texttt{tester} só acessa testes e especificações. & \texttt{tester} \\
    \hline
    PreToolUse & \texttt{Write|Edit} & Lista de proibições (\textit{deny-list}): o \texttt{dev} não altera arquivos de teste. & \texttt{dev} \\
    \hline
    SubagentStop & --- & Comparar os testes com o \textit{commit} da fase vermelha, executar a suíte e recusar o encerramento se algum teste foi alterado ou a suíte estiver vermelha. A implementação permanece [VERIFICAR]. & Global (\texttt{settings.json}) \\
    \hline
    PostToolUse & \texttt{*} & Registrar ferramenta, caminho e horário, sem o conteúdo, que pode conter dados pessoais, para auditoria do processo. & Global \\
    \hline
    PreToolUse (hipótese) & Chamada de subagente & Impedir o \texttt{tester} antes de a especificação ser esclarecida e o \texttt{dev} antes do \textit{commit} dos testes vermelhos. Hipótese de projeto, sem fonte; o nome da ferramenta que invoca subagentes deve ser conferido [VERIFICAR]. & Global \\
    \hline
    PreCompact (opcional) & --- & Salvar o estado da tarefa antes da compactação do contexto. & Global \\
    \hline
    \multicolumn{4}{p{0.95\textwidth}}{Fonte: autoria própria, com base em \citeonline{anthropic_hooks}.}\\
\end{longtable}
\endgroup

\subsubsection{Por que o Seed não utiliza MCP}
\label{subsubsec:mcp}

O \textit{Model Context Protocol} (MCP) conecta agentes a ferramentas e fontes de dados externas. O Seed não o adota por três razões de projeto: o desenvolvimento depende apenas das ferramentas nativas de leitura, busca, edição e execução e das skills descritas acima; cada servidor conectado amplia a superfície de acesso do agente, contrariando o princípio do privilégio mínimo \cite{saltzer_schroeder_1975}; e as descrições das ferramentas de cada servidor ocupam espaço no contexto [VERIFICAR: confirmar esse custo na documentação]. Essas razões são de desenho, não resultado de avaliação empírica.

% NOTA: a justificativa original mencionava o servidor "igris", que pertence a
% outro projeto (geração de currículo) e não deve aparecer na monografia.

\subsubsection{Isolamento do agente}
\label{subsubsec:isolamento}

O princípio do privilégio mínimo estabelece que cada componente deve acessar apenas o necessário para funcionar \cite{saltzer_schroeder_1975}. Para o agente, a camada principal é a execução em um contêiner ou ambiente isolado, sem acesso a variáveis sensíveis do sistema nem a diretórios fora do projeto. A documentação do Docker descreve, por exemplo, ambientes isolados em microVMs para agentes de código \cite{docker_sandboxes}. [VERIFICAR: citar a solução de isolamento de Akita, com repositório, e escolher a ferramenta efetivamente usada.] Os \textit{hooks} complementam esse isolamento, sem substituí-lo. Essa composição configura defesa em profundidade e integra os cuidados com os dados pessoais tratados pelo Seed no contexto da LGPD \cite{brasil2018lgpd}, sem constituir comprovação de conformidade legal. [VERIFICAR: conferir os artigos sobre dados pessoais sensíveis de saúde.]

%----------------------------------------------------------------------------
\subsection{Fluxo proposto}
\label{subsec:fluxo-ia}

O fluxo de cada feature compreende as etapas a seguir, ilustradas na Figura~\ref{fig:bpmn-ia}. Primeiro, o engenheiro descreve a feature e o orquestrador conduz o esclarecimento até fechar os critérios de aceitação e os limites numéricos. Em seguida, são geradas e revisadas a especificação, o planejamento e as tarefas. O \texttt{tester} produz os testes a partir da especificação, que são revisados, executados e registrados em um \textit{commit} na fase vermelha. O \texttt{dev} implementa a fase verde e refatora, mantendo os testes verdes. Por fim, a integração contínua e as ferramentas de segurança verificam o resultado, e o orquestrador reporta ao engenheiro.

% NOTA: incluir a figura (BPMN com raias Engenheiro, Orquestrador, tester, dev,
% hooks e ferramentas de segurança). Descomentar quando o arquivo existir.
% \begin{figure}[H]
%     \centering
%     \includegraphics[width=\textwidth]{editaveis/figuras/bpmn_ia.png}
%     \caption{BPMN do fluxo de desenvolvimento orientado por IA.}
%     \label{fig:bpmn-ia}
%     \par\small Fonte: autoria própria.
% \end{figure}

A verificação de segurança combina técnicas complementares: a análise estática (SAST) examina o código-fonte sem executá-lo; a análise dinâmica (DAST) testa a aplicação em execução; e a análise interativa (IAST) instrumenta a aplicação durante a execução para observar vulnerabilidades. [VERIFICAR: citar OWASP e o NIST SP 800-218 (SSDF) para essas definições e para a escolha das ferramentas.] A inclusão dessa etapa decorre das evidências de que código gerado por IA pode conter vulnerabilidades e de que o uso de assistentes aumenta a confiança indevida do desenvolvedor \cite{pearce_2021,perry_2022}. Um \textit{commit} que introduza vulnerabilidade crítica retorna à implementação.

%----------------------------------------------------------------------------
\subsection{Subagentes de teste e implementação}
\label{subsec:subagentes}

A divisão em dois subagentes busca a independência entre quem define o comportamento esperado e quem escreve o código. Seu fundamento é a observação de \citeonline{mathews_thesis_2025} e \citeonline{mathews_design_2024} de que ferramentas de geração de testes que inferem requisitos a partir do código tendem a validar defeitos. No estudo, as ferramentas avaliadas descartavam testes que falhavam no código com bugs e mantinham os que passavam, e em até 68,1\% das suítes finais de uma das ferramentas os testes passavam no código incorreto e falhavam no correto. Os autores defendem que o código seja escrito a partir de requisitos expressos como testes, de acordo com o TDD. Em um estudo preliminar com funções simples em Python, \citeonline{genc_2025} obtiveram testes executáveis e com cobertura total usando o código da função como entrada, mas o escopo era restrito a três funções e os autores reconhecem limitações com dependências entre componentes. Diferentemente dessas abordagens, o \texttt{tester} do Seed parte da especificação.

\subsubsection{Responsabilidades do subagente tester}

A descrição de \texttt{tester} atribui-lhe a produção de testes unitários a partir da especificação e dos critérios de aceitação de uma tarefa, após o esclarecimento de requisitos, denominado \textit{clarify}, e antes da implementação. Para operacionalizar essa atribuição, propõe-se que o orquestrador forneça a especificação da tarefa, os critérios de aceitação, os limites numéricos esclarecidos pelo engenheiro e o contrato de interfaces definido no planejamento, denominado \textit{plan}. Os limites devem vir do esclarecimento, e não ser inventados pelo modelo, pois bordas plausíveis nem sempre são as do domínio do Seed. A saída prevista consiste nos arquivos de teste e em uma síntese de sua localização e das pendências, sem reproduzir integralmente o código na conversa. A regra de devolver ambiguidades ao orquestrador, sem decidir requisitos, integra o protocolo proposto; sua presença nas instruções completas do agente permanece [VERIFICAR].

O Código~\ref{cod:tester} apresenta o trecho que seleciona o modelo, as ferramentas e a habilidade. O conjunto declarado contém \texttt{Read}, \texttt{Write}, \texttt{Edit}, \texttt{Grep} e \texttt{Glob}; \texttt{Bash} está deliberadamente ausente, conforme o comentário do arquivo. A exclusão do interpretador de comandos é coerente com a intenção de limitar os meios de acesso disponíveis para essa função, segundo o princípio do privilégio mínimo de \citeonline{saltzer_schroeder_1975}.

\begingroup
\setminted{firstline=4,lastline=7}
\codigoarq{autoria própria}{yaml}{editaveis/codigos/skills/tester.md}{Modelo, ferramentas e habilidade de tester}{cod:tester}
\endgroup

O mesmo arquivo associa o evento \texttt{PreToolUse} às cinco ferramentas por meio do padrão \texttt{Read|Grep|Glob|Write|Edit} e referencia o script \texttt{tester-allowlist.sh}. O nome do script, isoladamente, não comprova quais caminhos são autorizados. A política pretendida permite ler testes e especificações e escrever os testes, excluindo o código de produção. A opção por uma lista de permissões explícitas, ou \textit{allow-list}, decorre de permitir apenas recursos previamente identificados: uma lista de proibições, ou \textit{deny-list}, deixaria depender da enumeração de todas as localizações proibidas, e as buscas sem caminho varrem o repositório inteiro. Essa justificativa aplica o princípio de negação por padrão de \citeonline{saltzer_schroeder_1975}; a implementação concreta, incluindo o tratamento de caminhos e buscas, permanece [VERIFICAR].

A habilidade declarada é \texttt{unit-tests-boundary}. A documentação do Claude Code apresenta habilidades como instruções definidas em arquivos \texttt{SKILL.md} \cite{anthropic_skills}. Nas anotações do autor, sua descrição prevê JUnit~5 ou Jest, partição de equivalência e valores de borda. A partição organiza entradas em classes consideradas equivalentes para o objetivo de teste \cite{ammann_offutt_2008}; a análise de valores limite concentra casos nas fronteiras relevantes do domínio \cite{kuhn_coverage_2020}. Para o protocolo proposto, adota-se ainda a convenção AAA, com preparação do cenário (\textit{Arrange}), execução da operação (\textit{Act}) e verificação do resultado (\textit{Assert}). Essa convenção é uma escolha de organização dos testes; sua inclusão no corpo da habilidade permanece [VERIFICAR]. Os frameworks mantêm a escolha já registrada na Seção~\ref{sec:testes-software}, sem substituir a revisão dos resultados esperados pelo engenheiro.

\subsubsection{Responsabilidades do subagente dev}

A descrição de \texttt{dev} estabelece que sua execução ocorre depois da revisão e do registro dos testes vermelhos em um \textit{commit}. Cabe-lhe implementar o mínimo necessário para satisfazê-los e, em seguida, refatorar. Propõe-se que receba a especificação, o planejamento, a tarefa e a referência do commit de testes, acompanhados do resultado da execução que caracterizou a fase vermelha. A habilidade \texttt{implement-and-refactor}, declarada no arquivo, é descrita nas anotações como orientada às camadas Controller, Service, Model e Repository. O Código~\ref{cod:dev} evidencia que, além das ferramentas de leitura, busca e edição, esse agente dispõe de \texttt{Bash}.

\begingroup
\setminted{firstline=4,lastline=7}
\codigoarq{autoria própria}{yaml}{editaveis/codigos/skills/dev.md}{Modelo, ferramentas e habilidade de dev}{cod:dev}
\endgroup

No arquivo disponível, o hook \texttt{PreToolUse} de \texttt{dev} seleciona somente \texttt{Write|Edit} e referencia \texttt{dev-no-test-edit.sh}. A política pretendida é permitir alterações no código de produção e impedir alterações nos testes por essas ferramentas. Entretanto, os padrões da lista de proibições não puderam ser examinados [VERIFICAR]. Como \texttt{Bash} não integra esse seletor, a configuração apresentada não demonstra bloqueio de modificações de testes por comandos de terminal. A preocupação tem precedente: no agente de remediação de \citeonline{mathews_thesis_2025}, a instrução proíbe expressamente a modificação dos testes, o que indica que a alteração de testes é um desvio a ser evitado, embora o estudo não meça sua frequência. Caso um teste pareça contradizer a especificação, propõe-se interromper a implementação e devolver a divergência ao orquestrador, que consulta o engenheiro; o agente não deve ajustar o teste unilateralmente. A presença dessa regra na habilidade completa permanece [VERIFICAR].

Propõe-se também limitar as tentativas de correção. No estudo de \citeonline{mathews_thesis_2025}, o ciclo de remediação foi limitado a cinco iterações ou três falhas repetidas, e os autores observaram retorno decrescente além de três ou quatro tentativas. Esses parâmetros servem de ponto de partida para o Seed, a ser calibrado no TCC2.

\subsubsection{Encadeamento e pontos de controle}

O fluxo proposto é sequencial: o engenheiro esclarece os requisitos, o orquestrador entrega os insumos ao \texttt{tester}, e os testes produzidos são revisados e executados antes da delegação ao \texttt{dev}. Como o \texttt{tester} não dispõe de \texttt{Bash}, atribui-se a execução e o registro do commit ao engenheiro ou ao orquestrador com as permissões correspondentes. A revisão deve distinguir a falha esperada por funcionalidade ausente de problemas do ambiente ou de testes inválidos. O commit dos testes vermelhos constitui a referência de integridade para a implementação. Somente depois desse ponto se iniciam a fase verde e a refatoração, com nova execução dos testes, seguindo o ciclo discutido na Subseção~\ref{subsec:tdd} \cite{aniche_2014}. A fase verde e a refatoração não são paralelas: a refatoração só ocorre com os testes verdes. Não se prevê execução paralela dos dois agentes sobre a mesma tarefa.

A Tabela~\ref{tab:comparacao-subagentes} sintetiza a divisão proposta e os campos comprovados pelas cópias disponíveis. A visibilidade do \texttt{tester} depende tanto do conteúdo encaminhado pelo orquestrador quanto da implementação da restrição de acesso: limitar ferramentas não retira informações já incluídas na entrada da tarefa.

\begingroup
\small
\setlength{\tabcolsep}{4pt}
\renewcommand{\arraystretch}{1.16}
\begin{longtable}{>{\raggedright\arraybackslash}p{0.18\textwidth}>{\raggedright\arraybackslash}p{0.37\textwidth}>{\raggedright\arraybackslash}p{0.37\textwidth}}
    \caption{Comparação entre os subagentes do fluxo proposto}\label{tab:comparacao-subagentes}\\
    \hline
    \textbf{Aspecto} & \textbf{tester} & \textbf{dev} \\
    \hline
    \endfirsthead
    \hline
    \textbf{Aspecto} & \textbf{tester} & \textbf{dev} \\
    \hline
    \endhead
    Papel & Produzir testes na fase vermelha. & Implementar a fase verde e refatorar. \\
    Entradas propostas & Especificação, critérios de aceitação, limites do clarify e interfaces do plan. & Especificação, plan, tarefa, testes revisados e commit da fase vermelha. \\
    Ferramentas declaradas & Read, Write, Edit, Grep e Glob. & Read, Write, Edit, Grep, Glob e Bash. \\
    Visibilidade pretendida & Testes e especificações, incluindo contratos de interfaces. & Testes, especificações e código de produção. \\
    Restrições propostas & Não consultar produção nem resolver ambiguidades de requisitos. & Não editar testes nem decidir divergências entre testes e especificação. \\
    Hooks declarados & PreToolUse sobre as cinco ferramentas; referência a tester-allowlist.sh. & PreToolUse sobre Write e Edit; referência a dev-no-test-edit.sh. \\
    Verificação final & Não declarada na cópia examinada. & SubagentStop não declarado na cópia; configuração externa [VERIFICAR]. \\
    Habilidade declarada & \texttt{unit-tests-boundary}. & \texttt{implement-and-refactor}. \\
    Modelo declarado & \texttt{sonnet}. & \texttt{sonnet}. \\
    \hline
    \multicolumn{3}{p{0.94\textwidth}}{Fonte: autoria própria, com base nos arquivos reproduzidos no Apêndice~\ref{apendice:agentes-ia} e no protocolo proposto nesta subseção.}\\
\end{longtable}
\endgroup

\subsubsection{Salvaguardas e limites dos mecanismos}

A documentação de hooks informa que \texttt{PreToolUse} ocorre antes da chamada de ferramenta e permite recusá-la. Já \texttt{SubagentStop} ocorre ao término da resposta do subagente e pode bloquear seu encerramento; não desfaz operações anteriores \cite{anthropic_hooks}. Em \texttt{settings.json}, esse evento pode selecionar o tipo de agente. No frontmatter de um subagente, a documentação apresenta \texttt{Stop}, convertido em \texttt{SubagentStop} durante a execução \cite{anthropic_subagents}. Essas possibilidades documentadas não comprovam a existência de uma auditoria de integridade no Seed.

Propõe-se verificar, no encerramento de \texttt{dev}, se os testes permanecem iguais à referência registrada na fase vermelha. Essa verificação permitiria detectar posteriormente alterações persistentes feitas por vias não interceptadas, incluindo \texttt{Bash}, mas não impediria sua ocorrência. Sua implementação e abrangência permanecem [VERIFICAR]. A comparação final também não demonstraria, por si só, a ausência de alterações transitórias restauradas antes da auditoria. Por isso, instruções textuais sobre papéis são complementadas, na proposta, por restrições de ferramentas, hooks e revisão humana. O rigor específico dos hooks constitui decisão de projeto a avaliar, não prática consolidada demonstrada pela literatura citada.

Essa composição é concebida como defesa em profundidade. No desenho proposto, o isolamento por contêiner deve constituir a camada principal de separação dos recursos do ambiente, com acesso limitado a arquivos, credenciais e rede; os hooks exercem função complementar. Não foi disponibilizada configuração que permita verificar esse isolamento [VERIFICAR]. Sua adoção deve integrar os cuidados com os dados pessoais tratados pelo Seed, no contexto da LGPD já abordado no capítulo \cite{brasil2018lgpd}, sem ser apresentada como comprovação de conformidade legal.

%----------------------------------------------------------------------------
\subsection{Limitações e ameaças à validade}
\label{subsec:limitacoes-ia}

As duas configurações utilizam o mesmo seletor de modelo, \texttt{sonnet}, conforme os Códigos~\ref{cod:tester} e~\ref{cod:dev}. Assim, a separação de tarefas não demonstra independência dos erros: admite-se, como ameaça à validade, que ambos reproduzam interpretações e vieses semelhantes. Além disso, o seletor não registra uma versão fixa do modelo. Propõe-se avaliar no TCC2 a escolha de modelo por agente, registrando a identificação efetiva utilizada e comparando a adequação dos testes e das implementações sob condições controladas. A atribuição de modelos mais leves a tarefas de teste não deve ser assumida como premissa, pois a literatura sugere que a qualidade dos testes é determinante para a qualidade do código resultante.

Outra limitação diz respeito à transferência de resultados acadêmicos. \citeonline{mathews_tdd_2024} utilizaram testes escritos por humanos nos experimentos, enquanto, no protocolo aqui proposto, um subagente produz os testes. A própria dissertação afirma não responder quando e como a IA deve escrever testes \cite{mathews_thesis_2025}. Portanto, os resultados desse estudo não podem ser extrapolados como evidência de benefício da dupla \texttt{tester}/\texttt{dev}. Os experimentos também se limitam a funções e programas em Python, a modelos específicos e a \textit{benchmarks} conhecidos, com risco de vazamento de dados de treinamento, conforme reconhecido pelos autores, e não envolvem projetos em Java com Spring. Há ainda um risco para a revisão humana: o estudo com usuários de Fakhoury et al., relatado na dissertação, indica que participantes tiveram dificuldade de distinguir testes errados sugeridos pelo sistema \cite{mathews_thesis_2025}. A efetividade da revisão dos critérios e dos oráculos de teste pelo engenheiro permanece uma etapa metodológica proposta, cuja efetividade também deve ser avaliada.

Por fim, cobertura estrutural não prova qualidade nem atendimento integral dos requisitos: entradas relevantes podem permanecer ausentes mesmo quando o código é exercitado \cite{kuhn_coverage_2020}. Nos experimentos de \citeonline{mathews_thesis_2025}, por exemplo, o primeiro teste já atingia cobertura total em 75\% dos problemas do HumanEval e 92,4\% dos do MBPP. Propõe-se incluir teste de mutação na avaliação do TCC2, medindo a capacidade da suíte de detectar alterações artificiais no programa. \citeonline{just_mutants_2014} encontraram associação entre detecção de mutantes e de defeitos reais, independentemente da cobertura, no conjunto estudado. Essa evidência fundamenta a métrica complementar, sem transformá-la em prova universal de qualidade.

%----------------------------------------------------------------------------
\subsection{Plano de avaliação do processo no TCC2}
\label{subsec:avaliacao-ia}

Para que as afirmações sobre qualidade sejam sustentadas por dados, propõe-se registrar, por feature, as seguintes medidas: cobertura de testes e índice de mutantes detectados, com os operadores empregados e o tratamento de mutantes equivalentes; vulnerabilidades identificadas pelas ferramentas de análise e sua gravidade; retrabalho, medido pelo número de ciclos de correção e de devoluções ao engenheiro; violações das restrições registradas pelos \textit{hooks}; e ambiguidades devolvidas ao orquestrador. A comparação entre modelos por agente, citada na subseção anterior, deve usar as mesmas tarefas e condições. Os resultados serão apresentados no capítulo de resultados, junto da tabela de rastreabilidade entre histórias de usuário, requisitos, tarefas e testes, e a ausência de um grupo de controle e o tamanho reduzido da equipe devem ser declarados como limitações.